% Luc Destombe --- licence CC BY-NC-SA 4.0
% https://creativecommons.org/licenses/by-nc-sa/4.0/deed.fr
% Fichier autonome : compiler avec pdflatex, deux fois (signets, nombre de pages).
\def\lyceeversion{1.0}
\documentclass[11pt,a4paper]{article}
\makeatletter
\newif\ifeleve
\elevefalse

\newif\iflycee@cm
\newif\iflycee@palatino

\RequirePackage[utf8]{inputenc}
\RequirePackage[T1]{fontenc}
\RequirePackage[french]{babel}
\RequirePackage{etoolbox}

\newcommand{\ShorthandsOff}{\@ifundefined{shorthandoff}{}{\shorthandoff{;:!?}}}
\newcommand{\ShorthandsOn}{\@ifundefined{shorthandon}{}{\shorthandon{;:!?}}}
\ShorthandsOff
\AtBeginDocument{\ShorthandsOn}
\AtBeginEnvironment{tikzpicture}{\ShorthandsOff}

\RequirePackage{geometry}
\RequirePackage{fancyhdr}
\RequirePackage{lastpage}
\RequirePackage{multicol}
\RequirePackage{array}
\RequirePackage{enumitem}
\RequirePackage{xcolor}
\RequirePackage{needspace}

\RequirePackage{amsmath,amssymb}
\iflycee@cm\else\RequirePackage{fourier}\fi
\DeclareMathAlphabet{\mathcal}{OMS}{cmsy}{m}{n}
\renewcommand{\ttdefault}{lmtt}
\begingroup\catcode`\_=\active \catcode`\^=\active
\global\let_\sb \global\let^\sp
\endgroup
\AtBeginDocument{\catcode`\_=12 \mathcode`\_="8000
  \catcode`\^=12 \mathcode`\^="8000 }
\newlength{\retraitcalculs}
\setlength{\retraitcalculs}{2em}

\RequirePackage{tikz}
\usetikzlibrary{arrows.meta,calc,babel,shapes.misc}
\RequirePackage[most]{tcolorbox}
\tcbuselibrary{skins,breakable}

\definecolor{coulDef}{HTML}{1F4E79}
\definecolor{coulProp}{HTML}{7030A0}
\definecolor{coulRem}{HTML}{C55A11}
\definecolor{coulEx}{HTML}{2E7D32}
\definecolor{coulExo}{HTML}{B03A2E}
\definecolor{coulPart}{HTML}{1F4E79}
\definecolor{coulMeth}{HTML}{1B6E4F}
\definecolor{coulCourbe}{HTML}{0B5ED7}
\definecolor{coulCourbeB}{HTML}{C00000}
\definecolor{coulCourbeC}{HTML}{2E7D32}
\definecolor{coulGrille}{HTML}{8C8C8C}
\definecolor{coulRep}{HTML}{1B6E4F}
\definecolor{coulBar}{HTML}{2E7D32}

\newcommand{\R}{\mathbb{R}}
\newcommand{\Rs}{\mathbb{R}^{*}}
\newcommand{\N}{\mathbb{N}}
\newcommand{\Ns}{\mathbb{N}^{*}}
\newcommand{\Z}{\mathbb{Z}}
\newcommand{\Q}{\mathbb{Q}}
\renewcommand{\le}{\leqslant}
\renewcommand{\ge}{\geqslant}
\renewcommand{\approx}{\simeq}
\newcommand{\vect}[1]{\overrightarrow{#1}}
\newcommand{\norme}[1]{\left\lVert #1 \right\rVert}
\newcommand{\intervalle}[2]{\left[#1\,;#2\right]}
\RequirePackage{cancel}
\renewcommand{\CancelColor}{\color{coulExo}}
\newcommand{\fois}[1]{\times{\color{coulExo}#1}}
\newcommand{\barre}[1]{\cancel{#1}}
\newlength{\lycee@eqh}\newlength{\lycee@eqd}
\newsavebox{\lycee@eqa}\newsavebox{\lycee@eqb}\newsavebox{\lycee@eqc}
\newcommand{\lycee@eqcol}[1]{\hbox{$\renewcommand{\arraystretch}{1.35}%
  \begin{array}[t]{@{}>{\displaystyle}r@{}>{\displaystyle{}}l@{}}#1\end{array}$}}
\newcommand{\lycee@eqmesure}[1]{%
  \ifdim\ht#1>\lycee@eqh\setlength{\lycee@eqh}{\ht#1}\fi
  \ifdim\dp#1>\lycee@eqd\setlength{\lycee@eqd}{\dp#1}\fi}
\newcommand{\lycee@eqtrait}{\hspace{1.2em}\tikz[baseline=0pt]\draw[dash pattern=on 2pt off 2pt,
  line width=0.5pt] (0,\lycee@eqh) -- (0,-\lycee@eqd);\hspace{1.2em}}
\newcommand{\eqproduit}[2]{\par\smallskip\noindent C'est une équation produit :
  $A\times B=0 \iff A=0$ ou $B=0$.\par\nopagebreak\smallskip
  \begingroup\setlength{\lycee@eqh}{0pt}\setlength{\lycee@eqd}{0pt}%
  \setbox\lycee@eqa=\lycee@eqcol{#1}\setbox\lycee@eqb=\lycee@eqcol{#2}%
  \lycee@eqmesure\lycee@eqa\lycee@eqmesure\lycee@eqb%
  \noindent\hspace*{\retraitcalculs}\usebox\lycee@eqa\lycee@eqtrait\usebox\lycee@eqb\endgroup\par\smallskip}
\newcommand{\eqproduittrois}[3]{\par\smallskip\noindent C'est une équation produit :
  $A\times B\times C=0 \iff A=0$ ou $B=0$ ou $C=0$.\par\nopagebreak\smallskip
  \begingroup\setlength{\lycee@eqh}{0pt}\setlength{\lycee@eqd}{0pt}%
  \setbox\lycee@eqa=\lycee@eqcol{#1}\setbox\lycee@eqb=\lycee@eqcol{#2}\setbox\lycee@eqc=\lycee@eqcol{#3}%
  \lycee@eqmesure\lycee@eqa\lycee@eqmesure\lycee@eqb\lycee@eqmesure\lycee@eqc%
  \noindent\hspace*{\retraitcalculs}\usebox\lycee@eqa\lycee@eqtrait\usebox\lycee@eqb\lycee@eqtrait\usebox\lycee@eqc\endgroup
  \par\smallskip}

\tikzset{grille/.style={coulGrille,line width=0.4pt}}
\newcommand{\quadrillage}[4]{\draw[grille,step=1] (#1,#2) grid (#3,#4);}
\tikzset{
  croix/.style={cross out,draw,line width=0.6pt,inner sep=0pt,minimum size=4pt},
  croixdroite/.style={croix,rotate=45,line width=0.8pt},
}
\newcommand{\pt}[3]{\path (#1) node[croix]{} node[#2,font=\small,inner sep=2.5pt]{$#3$};}

\newcommand{\lycee@repere}[4]{%
  \draw[grille,step=1] (#1,#2) grid (#3,#4);
  \draw[-{Stealth[length=2mm]},thick] (#1,0)--({#3+0.4},0) node[below left=-1pt]{$x$};
  \draw[-{Stealth[length=2mm]},thick] (0,#2)--(0,{#4+0.4}) node[below left=-1pt]{$y$};
  \node[below left,font=\scriptsize,inner sep=1.5pt] at (0,0) {$0$};}
\newcommand{\lycee@gradx}[1]{\foreach \k/\l in {#1}{%
  \draw (\k,0.07)--(\k,-0.07) node[below,font=\scriptsize,inner sep=1.5pt]{$\l$};}}
\newcommand{\lycee@grady}[1]{\foreach \k/\l in {#1}{%
  \draw (0.07,\k)--(-0.07,\k) node[left,font=\scriptsize,inner sep=1.5pt]{$\l$};}}
\newcommand{\lycee@intff}[2]{\left[#1\,;#2\right]}
\newcommand{\lycee@intfo}[2]{\left[#1\,;#2\right[}
\newcommand{\lycee@intof}[2]{\left]#1\,;#2\right]}
\newcommand{\lycee@intoo}[2]{\left]#1\,;#2\right[}
\newcommand{\lycee@ens}[1]{\left\{#1\right\}}
\AtBeginDocument{%
  \providecommand{\repere}{\lycee@repere}%
  \providecommand{\gradx}{\lycee@gradx}%
  \providecommand{\grady}{\lycee@grady}%
  \providecommand{\Cf}{\mathcal{C}\sb{f}}%
  \providecommand{\Cg}{\mathcal{C}\sb{g}}%
  \providecommand{\intff}{\lycee@intff}%
  \providecommand{\intfo}{\lycee@intfo}%
  \providecommand{\intof}{\lycee@intof}%
  \providecommand{\intoo}{\lycee@intoo}%
  \providecommand{\ens}{\lycee@ens}}

\newlist{questions}{enumerate}{3}
\setlist[questions,1]{label=\arabic*\textdegree),leftmargin=*,itemsep=2pt,topsep=3pt}
\setlist[questions,2]{label=\alph*),leftmargin=*,itemsep=1pt,topsep=2pt}
\setlist[questions,3]{label=\roman*),leftmargin=*,itemsep=1pt,topsep=1pt}
\setlist[itemize]{label=\textbullet,leftmargin=*,itemsep=2pt,topsep=3pt}

\newcommand{\besoinplace}[1]{\par\penalty\@M\begingroup
  \dimen@=#1\relax\dimen@ii=\pagegoal\advance\dimen@ii-\pagetotal
  \ifdim\dimen@>\dimen@ii\ifdim\dimen@ii>\z@\vfil\fi\break\fi\endgroup}

\newcommand{\lycee@pied}{\thepage/\pageref{LastPage}}
\newcommand{\entete}[2]{%
  \pagestyle{fancy}\fancyhf{}%
  \fancyhead[L]{\footnotesize\color{coulPart}\textbf{#1}}%
  \fancyhead[R]{\footnotesize\color{coulPart}\textbf{#2}}%
  \fancyfoot[C]{\footnotesize\lycee@pied}%
  \renewcommand{\headrulewidth}{0.4pt}%
  \renewcommand{\headrule}{\hbox to\headwidth{%
    \color{coulPart!40}\leaders\hrule height \headrulewidth\hfill}}}

\RequirePackage{ccicons}
\newcommand{\lycee@licence}{{\scriptsize\color{gray}%
  \copyright\ Luc Destombe --- \ccbyncsa\ CC BY-NC-SA 4.0}}
\newif\iflycee@licence \lycee@licencetrue
\newcommand{\sanslicence}{\lycee@licencefalse}
\AtBeginDocument{\iflycee@licence\AddToHookNext{shipout/foreground}{%
  \put(0,0){\rlap{\kern\dimexpr1in+\hoffset+\oddsidemargin+\textwidth\relax
    \llap{\raisebox{-\dimexpr1in+\voffset+\topmargin+\headheight+\headsep
      +\textheight+\footskip\relax}{\lycee@licence}}}}}\fi}

\newcommand{\lycee@site}{faire-des-maths.fr}
\newcommand{\lycee@version}{{\scriptsize\color{gray}\lycee@site
  \ifdefined\lyceeversion\ --- Version~\lyceeversion\fi}}
\AtBeginDocument{\iflycee@licence\AddToHookNext{shipout/foreground}{%
  \put(0,0){\rlap{\kern\dimexpr1in+\hoffset+\oddsidemargin\relax
    \raisebox{-\dimexpr1in+\voffset+\topmargin+\headheight+\headsep
      +\textheight+\footskip\relax}{\lycee@version}}}}\fi}

\newcommand{\formulesserrees}{%
  \setlength{\abovedisplayskip}{3pt}\setlength{\belowdisplayskip}{3pt}%
  \setlength{\abovedisplayshortskip}{2pt}\setlength{\belowdisplayshortskip}{3pt}}

\setlength{\parindent}{0pt}

\RequirePackage[implicit=false,hidelinks,pdfencoding=auto,psdextra,bookmarksopen,
  bookmarksopenlevel=1,pdfcreator={},pdfproducer={}]{hyperref}
\RequirePackage{bookmark}
\hypersetup{pdfauthor={Luc Destombe},pdfkeywords={CC BY-NC-SA 4.0}}
\newcounter{lycee@signet}
\newcommand{\signet}[3][]{\stepcounter{lycee@signet}%
  \edef\lycee@dest{\if\relax\detokenize{#1}\relax
    signet.\arabic{lycee@signet}\else#1\fi}%
  \hypertarget{\lycee@dest}{}\bookmark[level=#2,dest=\lycee@dest]{#3}}
\pdfstringdefDisableCommands{%
  \def\R{R}\def\N{N}\def\Z{Z}\def\Q{Q}\def\vect#1{#1}\def\Cf{Cf}\def\Cg{Cg}\def\fois#1{×#1}%
  \def\barre#1{#1}\def\intervalle#1#2{[#1;#2]}\def\intff#1#2{[#1;#2]}%
  \def\textsuperscript#1{#1}\def\dfrac#1#2{#1/#2}\def\frac#1#2{#1/#2}%
  \def\vec#1{#1⃗}}%
\newcommand{\res}[1]{#1}
\geometry{top=2cm,bottom=1cm,left=1cm,right=1cm,headheight=14pt,footskip=0.6cm}

\RequirePackage{pgfplots}
\pgfplotsset{compat=1.18}
\RequirePackage{tkz-tab}

\newcounter{partiecours}
\renewcommand{\thepartiecours}{\Roman{partiecours}}
\newcommand{\partiecours}[1]{%
  \refstepcounter{partiecours}%
  \Needspace*{8\baselineskip}%
  \bigskip
  \noindent\begin{tcolorbox}[enhanced,colback=coulPart,colframe=coulPart,
      boxrule=0pt,arc=2pt,left=8pt,right=8pt,top=4pt,bottom=4pt,
      before skip=6pt,after skip=10pt]
    \color{white}\bfseries\large \thepartiecours{} --- #1%
    \signet[partie-\arabic{partiecours}]{1}{\thepartiecours{} --- #1}
  \end{tcolorbox}%
  \addcontentsline{toc}{section}{\thepartiecours{} --- #1}%
}

\newtcolorbox{boitecours}[3][]{%
  enhanced,breakable,
  colback=white,colframe=#2,coltitle=white,
  fonttitle=\bfseries,
  attach boxed title to top left={xshift=8pt,yshift=-\tcboxedtitleheight/2},
  boxed title style={colback=#2,arc=2pt,boxrule=0pt},
  boxrule=0.9pt,arc=2pt,
  left=8pt,right=8pt,top=8pt,bottom=6pt,
  title={#3},#1}

\newcounter{methode}
\newenvironment{definitionbox}[1][Définition]%
  {\begin{boitecours}[phantom={\signet{2}{#1}}]{coulDef}{#1}}{\end{boitecours}}
\newenvironment{proprietebox}[1][Propriété]%
  {\begin{boitecours}[phantom={\signet{2}{#1}}]{coulProp}{#1}}{\end{boitecours}}
\newenvironment{methodebox}[1][Méthode]{\stepcounter{methode}%
  \begin{boitecours}[phantom={\signet[methode-\arabic{methode}]{2}{#1}}]{coulMeth}{#1}}%
  {\end{boitecours}}

\newcommand{\etape}[1]{\par\smallskip\textbf{\color{coulMeth}#1}\par\nobreak\smallskip}
\newcommand{\enonce}[1]{\emph{#1}\par\smallskip}
\newcommand{\reponse}[1]{\par\smallskip
  {\footnotesize\itshape\color{black!55}Réponses : #1}}
\newcommand{\demo}[1]{\textbf{\color{coulMeth}Démonstration.} #1}
\newcommand{\afaire}[1]{%
  \par\smallskip
  \noindent{\small\color{coulExo}$\blacktriangleright$\ \textbf{À faire :}
    fiche d'exercices du chapitre, #1.}\par\smallskip}

\newenvironment{calculs}{\setlength{\jot}{5pt}\@fleqntrue\@mathmargin\retraitcalculs
  \setlength{\abovedisplayskip}{4pt}\setlength{\belowdisplayskip}{4pt}%
  \setlength{\abovedisplayshortskip}{4pt}\setlength{\belowdisplayshortskip}{4pt}%
  \csname alignat*\endcsname{2}}%
  {\csname endalignat*\endcsname}
\newcommand{\rem}[1]{\text{\small\itshape\color{black!60}#1}}
\newcommand{\explique}[2]{\par\smallskip\noindent
  \begin{minipage}[t]{0.57\linewidth}\raggedright #1\end{minipage}\hfill
  \begin{minipage}[t]{0.40\linewidth}\raggedright\leavevmode
    {\small\itshape\color{black!60}#2}\end{minipage}\par}
\newcommand{\cas}[1]{\par\smallskip\noindent\textbf{\color{coulExo}#1}\vadjust{\nobreak}}
\newcommand{\calc}[1]{\par\smallskip\textbf{\color{coulExo}#1}\par\nobreak}

\newtcolorbox{remarquebox}[1][Remarque]{%
  enhanced,breakable,colback=white,colframe=coulRem,
  boxrule=0pt,leftrule=2.5pt,arc=0pt,outer arc=0pt,
  left=8pt,right=6pt,top=5pt,bottom=5pt,
  before upper={\textbf{\color{coulRem}#1}\par\smallskip}}

\newtcolorbox{exemplebox}[1][Exemple]{%
  enhanced,breakable,colback=white,colframe=coulEx,
  boxrule=0pt,leftrule=2.5pt,arc=0pt,outer arc=0pt,
  left=8pt,right=6pt,top=5pt,bottom=5pt,
  before upper={\textbf{\color{coulEx}#1}\par\smallskip}}

\pgfplotsset{
  repere/.style={
    axis lines=middle,
    axis line style={-{Stealth[length=2.4mm]},thick},
    every axis x label/.style={at={(ticklabel* cs:1.0)},anchor=west},
    every axis y label/.style={at={(ticklabel* cs:1.0)},anchor=south},
    label style={font=\footnotesize},
    tick label style={font=\scriptsize},
    grid=both,
    major grid style={grille},
    minor tick num=0,
    tick align=inside,
    clip mode=individual,
  },
}
\tikzset{
  courbe/.style={coulCourbe,line width=1.1pt,smooth},
  courbeB/.style={coulCourbeB,line width=1.1pt,smooth},
  courbeC/.style={coulCourbeC,line width=1.1pt,smooth},
}

\newlist{etapes}{enumerate}{1}
\setlist[etapes]{label=\textbf{\arabic*.},leftmargin=*,itemsep=1pt,topsep=2pt}

\newlist{expressions}{enumerate}{1}
\setlist[expressions]{label={\Alph*\ $=$},leftmargin=*,itemsep=3pt,topsep=3pt}

\newcounter{exo}
\renewcommand{\theexo}{\ifnum\value{exo}<10 0\fi\arabic{exo}}
\newtcolorbox{exercice}[1][]{%
  enhanced,breakable,
  colback=white,colframe=coulExo,
  boxrule=0.9pt,arc=2pt,
  left=8pt,right=8pt,top=8pt,bottom=6pt,
  coltitle=white,fonttitle=\bfseries,
  attach boxed title to top left={xshift=8pt,yshift=-\tcboxedtitleheight/2},
  boxed title style={colback=coulExo,arc=2pt,boxrule=0pt},
  phantom={\signet{2}{Exercice \arabic{exo}}},
  title={Exercice \theexo\ifx\relax#1\relax\else{} --- #1\fi}}
\newenvironment{exo}[1][\relax]{\refstepcounter{exo}\begin{exercice}[#1]}{\end{exercice}}

  \newenvironment{correction}{\par}{\par}
  \newcommand{\autableau}{}
  \newcommand{\etapeexemples}{\etape{Exemples corrigés}}
  \newcommand{\etapetableau}[1]{\etape{#1}}
  \newenvironment{exempletableau}[1][Exemple]%
    {\begin{exemplebox}[#1]}{\end{exemplebox}}

\RequirePackage{comment}
\excludecomment{correctionavj}

\newcommand{\coursEntete}{}
\newcommand{\coursNiveau}{}
\newcommand{\coursTitre}{}
\newcommand{\coursSurTitre}{}
\newcommand{\coursSousTitre}{}

\newcommand{\enteteCours}[1]{\renewcommand{\coursEntete}{#1}}
\newcommand{\chapitrecours}[2]{\enteteCours{Chapitre #1 --- #2}}
\newcommand{\niveaucours}[1]{\renewcommand{\coursNiveau}{#1}}
\newcommand{\titrecours}[3][]{%
  \renewcommand{\coursTitre}{#2}%
  \renewcommand{\coursSurTitre}{#1}%
  \renewcommand{\coursSousTitre}{#3}}

\pagestyle{fancy}
\fancyhf{}
\fancyhead[L]{\footnotesize\color{coulPart}\textbf{\coursEntete}}
\fancyhead[R]{\footnotesize\color{coulPart}\textbf{\coursNiveau}}
\fancyfoot[C]{\footnotesize\thepage}
\renewcommand{\headrulewidth}{0.4pt}
\renewcommand{\headrule}{\hbox to\headwidth{\color{coulPart!40}\leaders\hrule height \headrulewidth\hfill}}

\setlength{\parskip}{2pt}

\newcommand{\lycee@mentioneleve}{\\[2pt]{\small\itshape Version élève}}
\AtBeginDocument{%
  \begin{center}
    \begin{tcolorbox}[enhanced,width=0.72\linewidth,colback=white,colframe=coulPart,
        boxrule=1.6pt,arc=3pt,halign=center,top=10pt,bottom=10pt]
      \ifx\coursSurTitre\empty
        {\Huge\bfseries\color{coulPart} \coursTitre}\\[4pt]
      \else
        {\Huge\bfseries\color{coulPart} \coursTitre}\\[3pt]
        {\Large\bfseries\color{coulPart} \coursSurTitre}\\[5pt]
      \fi
      {\normalsize \coursSousTitre}%
      
    \end{tcolorbox}
  \end{center}%
}

\RequirePackage{tikz}
\RequirePackage{xcolor}

\usetikzlibrary{decorations.text}
\definecolor{lycee@metal}{HTML}{1F4E79}
\definecolor{lycee@manche}{HTML}{8A5A2B}
\definecolor{lycee@etincelle}{HTML}{F28C28}
\newcommand{\lycee@pluriel}[2]{#1 #2\ifnum#1>1 s\fi}
\newcommand{\lycee@sceau}{%
  \edef\lycee@bas{\lycee@pluriel{\lyceerelectures}{relecture} et
    \lycee@pluriel{\lyceerodages}{rodage}}%
  \begin{tikzpicture}[line join=round]
    \fill[coulPart!10,draw=coulPart,line width=.8pt] (0:2.2)
      \foreach \i in {1,...,71}{ -- (\i*5:{isodd(\i)?2.1:2.2}) } -- cycle;
    \draw[coulPart,line width=1.2pt] (0,0) circle (1.93);
    \draw[coulPart,line width=.4pt] (0,0) circle (1.85);
    \draw[coulPart,line width=.4pt] (0,0) circle (1.2);
    \path[decorate,decoration={text along path,text align=center,
        text={|\fontsize{9.5}{11}\selectfont\bfseries\color{coulPart}|TEST{\'E} AVEC MES {\'E}L{\`E}VES}}]
      (180:1.45) arc (180:0:1.45);
    \path[decorate,decoration={text along path,text align=center,
        text/.expanded={|\noexpand\fontsize{8.5}{10}\noexpand\selectfont\noexpand\color{coulPart}|\lycee@bas}}]
      (195:1.68) arc (195:345:1.68);
    \begin{scope}[x=.25mm,y=-.25mm]
      \begin{scope}[shift={(-32,-40)}]%
        \fill[lycee@metal] (23,31) rectangle (41,49);
        \draw[lycee@etincelle,line width=.5pt,line cap=round]
          (20,30)--(15,27) (19,34)--(14,35) (44,30)--(49,27) (45,34)--(50,35);
        \fill[lycee@manche,rounded corners=.3mm] (38,24.6) rectangle (60,27.6);
        \fill[lycee@metal,rounded corners=.4mm] (25,20) rectangle (39,31);
        \fill[lycee@etincelle,rounded corners=.2mm] (25,20) rectangle (28,31);
      \end{scope}
    \end{scope}
    \ifdefined\lyceesceauplusieurs\def\lycee@annee{Ann{\'e}es~:}\else
      \def\lycee@annee{Ann{\'e}e}\fi
    \path[decorate,decoration={text along path,text align=center,
        text/.expanded={|\noexpand\fontsize{8}{10}\noexpand\selectfont\noexpand\color{coulPart}|%
          \unexpanded\expandafter{\lycee@annee} \unexpanded\expandafter{\lyceesceauannees}}}]
      (95:1.08) arc (95:445:1.08);
  \end{tikzpicture}}
\AtEndDocument{\ifdefined\lyceesceauannees
  \par\besoinplace{5cm}\vspace{1.5em}%
  \begin{center}\lycee@sceau\end{center}\fi}
\makeatother

\chapitrecours{10}{Probabilités conditionnelles}
\niveaucours{1\textsuperscript{re} STI2D}
\titrecours{PROBABILITÉS CONDITIONNELLES}{Cours --- Première STI2D}

\begin{document}

\newcommand{\venn}[6]{%
  \draw (-1.75,-1.35) rectangle (2.75,1.35);
  \node[font=\footnotesize] at (-1.5,1.1) {$\Omega$};
  \draw[coulCourbe,thick] (0,0) circle (1.1);
  \draw[coulCourbeB,thick] (1,0) circle (1.1);
  \node[coulCourbe] at (-0.95,0.95) {$#1$};
  \node[coulCourbeB] at (1.95,0.95) {$#2$};
  \node at (-0.6,0) {#3};
  \node at (0.5,0) {#4};
  \node at (1.6,0) {#5};
  \node at (2.3,-1.05) {#6};}
\newcommand{\vennbord}{%
  \draw (-1.75,-1.35) rectangle (2.75,1.35);
  \draw (0,0) circle (1.1) (1,0) circle (1.1);
  \node[font=\scriptsize] at (-0.95,0.95) {$A$};
  \node[font=\scriptsize] at (1.95,0.95) {$B$};}
\tikzset{evt/.style={inner sep=2pt},
  proba/.style={font=\footnotesize,fill=white,inner sep=1pt},
  arbA/.style={},arbnA/.style={},arbAB/.style={},arbAnB/.style={},
  arbnAB/.style={},arbnAnB/.style={},chemin/.style={coulExo,line width=1.4pt}}
\newcommand{\arbre}[8]{%
  \node (O) at (0,0) {};
  \node[evt] (A) at (2.6,1.4) {$#1$};
  \node[evt] (nA) at (2.6,-1.4) {$\overline{#1}$};
  \node[evt] (AB) at (5.2,2.1) {$#2$};
  \node[evt] (AnB) at (5.2,0.7) {$\overline{#2}$};
  \node[evt] (nAB) at (5.2,-0.7) {$#2$};
  \node[evt] (nAnB) at (5.2,-2.1) {$\overline{#2}$};
  \draw[arbA] (O.center) -- node[proba]{$#3$} (A);
  \draw[arbnA] (O.center) -- node[proba]{$#4$} (nA);
  \draw[arbAB] (A) -- node[proba]{$#5$} (AB);
  \draw[arbAnB] (A) -- node[proba]{$#6$} (AnB);
  \draw[arbnAB] (nA) -- node[proba]{$#7$} (nAB);
  \draw[arbnAnB] (nA) -- node[proba]{$#8$} (nAnB);}
\newcommand{\feuille}[2]{\node[right,font=\footnotesize] at (#1.east) {#2};}
\newcommand{\feuilles}[2]{\feuille{AB}{$#1\cap #2$}\feuille{AnB}{$#1\cap\overline{#2}$}%
  \feuille{nAB}{$\overline{#1}\cap #2$}\feuille{nAnB}{$\overline{#1}\cap\overline{#2}$}}

\begin{remarquebox}[Comment utiliser ce cours]
Chaque partie commence par ce qu'il faut \textbf{savoir} (définitions, propriétés), puis
vient ce qu'il faut \textbf{savoir faire} : une méthode, avec des
\textbf{exemples corrigés} faits en classe, puis des exercices
\og\textbf{À vous de jouer}\fg{} à chercher seul.

\end{remarquebox}

\partiecours{Vocabulaire des événements}

On lance un dé : on connaît tous les résultats possibles, mais on ne peut pas prévoir
celui qui sortira. Les probabilités mesurent les chances de chacun.

\begin{definitionbox}[Définition 1 --- Expérience aléatoire, issue, univers]
Une expérience est \textbf{aléatoire} quand on connaît tous ses résultats possibles,
appelés \textbf{issues}, sans pouvoir prévoir celui qui se produira.

\smallskip
L'ensemble de toutes les issues est l'\textbf{univers}, noté $\Omega$ (\og oméga\fg{}).
Pour un lancer de dé à six faces : $\Omega=\ens{1\,;2\,;3\,;4\,;5\,;6}$.
\end{definitionbox}

\begin{definitionbox}[Définition 2 --- Événement]
Un \textbf{événement} est un ensemble d'issues. On le décrit par une phrase, ou par la
liste de ses issues entre accolades. Il est \textbf{réalisé} quand l'issue obtenue en
fait partie.

\smallskip
Par exemple, pour un lancer de dé, l'événement $A$ \og obtenir un nombre pair\fg{} s'écrit
$A=\ens{2\,;4\,;6}$ ; si on obtient $4$, l'événement $A$ est réalisé.

\smallskip
Un événement qui ne contient aucune issue est \textbf{impossible}, noté $\varnothing$ ;
l'univers $\Omega$ est l'événement \textbf{certain}.
\end{definitionbox}

\begin{definitionbox}[Définition 3 --- Intersection, réunion, contraire]
Soit $A$ et $B$ deux événements d'une même expérience.

\smallskip
\renewcommand{\arraystretch}{1.3}
\begin{tabular}{@{}>{\raggedright\arraybackslash}m{0.25\linewidth}
  >{\raggedright\arraybackslash}m{0.47\linewidth}>{\centering\arraybackslash}m{0.18\linewidth}@{}}
$A\cap B$ \newline \og $A$ inter $B$\fg{}, \og $A$ et $B$\fg{}
  & les issues qui sont \textbf{à la fois} dans $A$ et dans $B$
  & \begin{tikzpicture}[scale=0.42,baseline=(current bounding box.center)]
      \begin{scope}\clip (0,0) circle (1.1); \fill[coulMeth!30] (1,0) circle (1.1);\end{scope}
      \vennbord
    \end{tikzpicture}\\
$A\cup B$ \newline \og $A$ union $B$\fg{}, \og $A$ ou $B$\fg{}
  & les issues qui sont dans $A$, dans $B$ ou dans les deux : \textbf{au moins un}
    des deux événements est réalisé
  & \begin{tikzpicture}[scale=0.42,baseline=(current bounding box.center)]
      \fill[coulMeth!30] (0,0) circle (1.1) (1,0) circle (1.1);
      \vennbord
    \end{tikzpicture}\\
$\overline{A}$ \newline \og $A$ barre\fg{}, contraire de $A$
  & les issues de $\Omega$ qui \textbf{ne sont pas} dans $A$
  & \begin{tikzpicture}[scale=0.42,baseline=(current bounding box.center)]
      \fill[coulMeth!30,even odd rule] (-1.75,-1.35) rectangle (2.75,1.35) (0,0) circle (1.1);
      \vennbord
    \end{tikzpicture}\\
$A$ et $B$ \textbf{incompatibles}
  & $A\cap B=\varnothing$ : aucune issue commune ; $A$ et $B$ ne peuvent pas être
    réalisés en même temps
  & \begin{tikzpicture}[scale=0.42,baseline=(current bounding box.center)]
      \draw (-1.75,-1.35) rectangle (2.75,1.35);
      \draw (-0.3,0) circle (0.9) (1.6,0) circle (0.9);
      \node[font=\scriptsize] at (-1.25,0.95) {$A$};
      \node[font=\scriptsize] at (2.5,0.95) {$B$};
    \end{tikzpicture}
\end{tabular}
\end{definitionbox}

\begin{methodebox}[Méthode 1 --- Décrire des événements]
\etapeexemples
On lance un dé équilibré à $12$ faces, numérotées de $1$ à $12$. On note $A$ \og obtenir un
nombre pair\fg{}, $B$ \og obtenir un multiple de $3$\fg{} et $C$ \og obtenir un nombre impair
inférieur à $6$\fg{}.
\begin{questions}[label=\alph*)]
  \item Écrire l'univers $\Omega$, puis les événements $A$ et $B$ sous forme d'ensembles.
  \item Décrire par une phrase, puis écrire sous forme d'ensembles, les événements $A\cap B$
        et $A\cup B$.
  \item Écrire l'événement $\overline{A}$ sous forme d'ensemble, puis le décrire par une phrase.
  \item Les événements $A$ et $C$ sont-ils incompatibles ? Et les événements $B$ et $C$ ?
\end{questions}
\begin{correction}
\cas{a)}
\explique{$\Omega=\res{\ens{1\,;2\,;3\,;\dots\,;12}}$ : $12$ issues.}{toutes les faces du dé}
\explique{$A=\res{\ens{2\,;4\,;6\,;8\,;10\,;12}}$}{les nombres pairs de $1$ à $12$}
\explique{$B=\res{\ens{3\,;6\,;9\,;12}}$}{les multiples de $3$ jusqu'à $12$}
\cas{b)}
\explique{$A\cap B$ : \og obtenir un nombre pair \textbf{et} multiple de $3$\fg{} ;
$A\cap B=\res{\ens{6\,;12}}$.}{les issues communes à $A$ et à $B$}
\explique{$A\cup B$ : \og obtenir un nombre pair \textbf{ou} un multiple de $3$\fg{} ;
$A\cup B=\res{\ens{2\,;3\,;4\,;6\,;8\,;9\,;10\,;12}}$.}{les issues de $A$, puis celles de $B$
qui manquent : $3$ et $9$}
\cas{c)}
\explique{$\overline{A}=\res{\ens{1\,;3\,;5\,;7\,;9\,;11}}$ : \og obtenir un nombre
\res{impair}\fg{}.}{les issues de $\Omega$ qui ne sont pas dans $A$}
\cas{d)}
\explique{$C=\ens{1\,;3\,;5}$ et $A\cap C=\varnothing$, donc \res{$A$ et $C$ sont
incompatibles}.}{aucun nombre n'est à la fois pair et impair}
\explique{$B\cap C=\ens{3}$, donc \res{$B$ et $C$ ne sont pas incompatibles}.}{si on
obtient $3$, $B$ et $C$ sont réalisés en même temps}
\end{correction}

\etape{À vous de jouer}
On lance un dé équilibré à $8$ faces, numérotées de $1$ à $8$. On note $A$ \og obtenir un
nombre impair\fg{}, $B$ \og obtenir un nombre supérieur ou égal à $5$\fg{} et $C$ \og obtenir un
multiple de $4$\fg{}.
\begin{questions}[label=\alph*)]
  \item Écrire l'univers $\Omega$, puis les événements $A$ et $B$ sous forme d'ensembles.
  \item Décrire par une phrase, puis écrire sous forme d'ensembles, les événements $A\cap B$
        et $A\cup B$.
  \item Écrire l'événement $\overline{A}$ sous forme d'ensemble, puis le décrire par une phrase.
  \item Les événements $A$ et $C$ sont-ils incompatibles ? Et les événements $B$ et $C$ ?
\end{questions}
\reponse{a) $\Omega=\ens{1\,;2\,;\dots\,;8}$ ; $A=\ens{1\,;3\,;5\,;7}$ ;
$B=\ens{5\,;6\,;7\,;8}$ ;\; b)~$A\cap B=\ens{5\,;7}$ ;
$A\cup B=\ens{1\,;3\,;5\,;6\,;7\,;8}$ ;\; c)~$\overline{A}=\ens{2\,;4\,;6\,;8}$, \og obtenir
un nombre pair\fg{} ;\; d)~$A$ et $C$ incompatibles ; $B$ et $C$ non ($B\cap C=\ens{8}$).}
\end{methodebox}

\afaire{exercices 1 à 4}

\partiecours{Diagramme de Venn}

Dans une population (les membres d'un club, les emplacements d'un camping), on range
chaque élément selon deux critères. Un diagramme de Venn montre combien d'éléments
vérifient l'un, l'autre, les deux ou aucun.

\begin{definitionbox}[Définition 4 --- Diagramme de Venn]
\begin{minipage}[c]{0.56\linewidth}\raggedright
Un \textbf{diagramme de Venn} représente l'univers par un rectangle et deux événements
$A$ et $B$ par deux cercles. Il partage la population en quatre zones, sans
chevauchement :
\begin{itemize}
  \item $A\cap B$ : dans les deux cercles ;
  \item $A\cap\overline{B}$ : dans $A$ seulement ;
  \item $\overline{A}\cap B$ : dans $B$ seulement ;
  \item $\overline{A}\cap\overline{B}$ : hors des deux cercles (ni $A$ ni $B$).
\end{itemize}
\end{minipage}\hfill
\begin{minipage}[c]{0.4\linewidth}\centering
\begin{tikzpicture}[scale=0.95]
  \venn{A}{B}{\footnotesize$A\cap\overline{B}$}{\footnotesize$A\cap B$}%
    {\footnotesize$\overline{A}\cap B$}{\footnotesize$\overline{A}\cap\overline{B}$}
\end{tikzpicture}
\end{minipage}
\end{definitionbox}

\begin{remarquebox}[À retenir]
L'effectif de $A$ est la somme des \textbf{deux} zones de son cercle : \og dans $A$\fg{} n'est
pas \og dans $A$ seulement\fg{}. On remplit donc le diagramme en commençant par la zone
commune.
\end{remarquebox}

\begin{methodebox}[Méthode 2 --- Compléter un diagramme de Venn]
\etapeexemples
Un club de robotique compte $40$ membres : $18$ programment les robots (ensemble $I$,
informatique), $25$ conçoivent les pièces (ensemble $M$, mécanique) et $9$ font les deux.
\begin{questions}[label=\alph*)]
  \item Représenter la situation par un diagramme de Venn.
  \item Combien de membres font au moins une des deux activités ? Aucune ? Une seule ?
\end{questions}
Un camping compte $50$ emplacements ; on note $E$ l'ensemble de ceux qui ont
l'électricité et $O$ l'ensemble de ceux qui ont un point d'eau. $30$ emplacements ont
l'électricité, $5$ ont l'électricité et un point d'eau, $10$ n'ont ni l'un ni l'autre.
\begin{questions}[label=\alph*),start=3]
  \item Représenter la situation par un diagramme de Venn.
  \item Combien d'emplacements ont un point d'eau ?
\end{questions}
\begin{correction}
\cas{a)}\par\nopagebreak
\begin{minipage}[c]{0.58\linewidth}\raggedright
On part de la zone commune : $9$ membres font les deux.
\begin{calculs}
18-9 &= \res{9} &\qquad& \rem{informatique seulement}\\
25-9 &= \res{16} && \rem{mécanique seulement}\\
40-(9+9+16) &= \res{6} && \rem{ni l'un ni l'autre}
\end{calculs}
\end{minipage}\hfill
\begin{minipage}[c]{0.38\linewidth}\centering
\begin{tikzpicture}[scale=0.8]
  \venn{I}{M}{$9$}{$9$}{$16$}{$6$}
\end{tikzpicture}
\end{minipage}
\cas{b)}
\begin{calculs}
9+9+16 &= \res{34} &\qquad& \rem{au moins une activité : les trois zones des cercles}\\
9+16 &= \res{25} && \rem{une seule : les deux zones \og seulement\fg{}}
\end{calculs}
\explique{$34$ membres font au moins une activité, \res{$6$} n'en font aucune et $25$ en
font une seule.}{aucune : la zone hors des cercles}
\cas{c)}\par\nopagebreak
\begin{minipage}[c]{0.58\linewidth}\raggedright
On part de la zone commune : $5$ emplacements ont les deux.
\begin{calculs}
30-5 &= \res{25} &\qquad& \rem{électricité seulement}\\
50-10 &= 40 && \rem{dans au moins un des cercles}\\
40-(25+5) &= \res{10} && \rem{point d'eau seulement}
\end{calculs}
\end{minipage}\hfill
\begin{minipage}[c]{0.38\linewidth}\centering
\begin{tikzpicture}[scale=0.8]
  \venn{E}{O}{$25$}{$5$}{$10$}{$10$}
\end{tikzpicture}
\end{minipage}
\cas{d)}
\begin{calculs}
5+10 &= \res{15} &\qquad& \rem{les deux zones du cercle $O$}
\end{calculs}
\explique{$15$ emplacements ont un point d'eau.}{\og point d'eau\fg{} n'est pas \og point d'eau
seulement\fg{}}
\end{correction}

\etape{À vous de jouer}
Lors des portes ouvertes d'un lycée, $120$ visiteurs sont venus : $70$ ont assisté à la
démonstration de drones (ensemble $D$), $50$ à l'atelier d'impression 3D (ensemble $I$) et
$30$ aux deux.
\begin{questions}[label=\alph*)]
  \item Représenter la situation par un diagramme de Venn.
  \item Combien de visiteurs ont assisté à au moins une des deux activités ? À aucune ? À
        une seule ?
\end{questions}
Une résidence compte $80$ logements ; on note $S$ l'ensemble de ceux qui ont des panneaux
solaires et $C$ l'ensemble de ceux qui ont une pompe à chaleur. $45$ logements ont des
panneaux solaires, $15$ ont des panneaux solaires et une pompe à chaleur, $20$ n'ont ni l'un
ni l'autre.
\begin{questions}[label=\alph*),start=3]
  \item Représenter la situation par un diagramme de Venn.
  \item Combien de logements ont une pompe à chaleur ?
\end{questions}
\reponse{a) $D$ seulement $40$, les deux $30$, $I$ seulement $20$, aucun $30$ ;\;
b)~$90$ ; $30$ ; $60$ ;\; c)~$S$ seulement $30$, les deux $15$, $C$ seulement $15$, aucun
$20$ ;\; d)~$30$.}
\end{methodebox}

\afaire{exercices 5 à 8}

\partiecours{Équiprobabilité}

\begin{definitionbox}[Définition 5 --- Probabilité d'un événement]
La \textbf{probabilité} d'un événement $A$, notée $P(A)$, est un nombre compris entre $0$
et $1$ qui mesure la chance que $A$ a de se réaliser : $P(A)=0{,}8$ signifie que $A$ a
$8$ chances sur $10$ de se réaliser, soit $80\,\%$.

\smallskip
Un événement impossible a une probabilité égale à $0$, l'événement certain une
probabilité égale à $1$.
\end{definitionbox}

\begin{proprietebox}[Propriété 1 --- Équiprobabilité]
Quand toutes les issues ont la même probabilité (dé équilibré, tirage au hasard, boules
indiscernables au toucher), on dit qu'il y a \textbf{équiprobabilité}, et
\[
  P(A)=\dfrac{\text{nombre d'issues de $A$}}{\text{nombre d'issues de $\Omega$}}.
\]
Quand on choisit au hasard un élément d'une population, $P(A)$ est l'effectif de $A$
divisé par l'effectif total.
\end{proprietebox}

\begin{methodebox}[Méthode 3 --- Calculer une probabilité dans une situation d'équiprobabilité]
\etapeexemples
\begin{minipage}[c]{0.6\linewidth}\raggedright
\begin{questions}[label=\alph*)]
  \item On lance un dé équilibré à $12$ faces, numérotées de $1$ à $12$. On note $A$
        \og obtenir un nombre pair\fg{} et $B$ \og obtenir un multiple de $3$\fg{}. Calculer $P(A)$,
        $P(B)$ et $P(A\cap B)$.
  \item Le diagramme ci-contre décrit les $40$ membres d'un club de robotique ($I$ :
        informatique ; $M$ : mécanique). On choisit un membre au hasard. Calculer
        $P(I\cap M)$, $P(I\cup M)$ et $P(\overline{I}\cap\overline{M})$.
\end{questions}
\end{minipage}\hfill
\begin{minipage}[c]{0.36\linewidth}\centering
\begin{tikzpicture}[scale=0.8]
  \venn{I}{M}{$9$}{$9$}{$16$}{$6$}
\end{tikzpicture}
\end{minipage}
\begin{correction}
\cas{a)}
\explique{Le dé est équilibré : il y a équiprobabilité, avec $12$ issues.}{chaque face a
la même chance}
\begin{calculs}
P(A) &= \dfrac{6}{12} &\qquad& \rem{$A=\ens{2\,;4\,;6\,;8\,;10\,;12}$ ; $6$ issues sur $12$}\\
P(A) &= \dfrac{\barre{6}\times 1}{\barre{6}\times 2} && \rem{on simplifie par $6$}\\
P(A) &= \res{\dfrac{1}{2}} &&\\
P(B) &= \dfrac{4}{12} = \res{\dfrac{1}{3}} && \rem{$B=\ens{3\,;6\,;9\,;12}$ ; on
simplifie par $4$}\\
P(A\cap B) &= \dfrac{2}{12} = \res{\dfrac{1}{6}} && \rem{$A\cap B=\ens{6\,;12}$ ; on
simplifie par $2$}
\end{calculs}
\cas{b)}
\explique{Le membre est choisi au hasard : il y a équiprobabilité, et on divise chaque
effectif par $40$.}{$40$ membres en tout}
\begin{calculs}
P(I\cap M) &= \dfrac{9}{40} = \res{0{,}225} &\qquad& \rem{la zone commune}\\
P(I\cup M) &= \dfrac{9+9+16}{40} && \rem{les trois zones des cercles}\\
P(I\cup M) &= \dfrac{34}{40} = \res{0{,}85} &&\\
P(\overline{I}\cap\overline{M}) &= \dfrac{6}{40} = \res{0{,}15} && \rem{ni informatique ni
mécanique}
\end{calculs}
\explique{Un membre choisi au hasard fait au moins une des deux activités avec une
probabilité de $0{,}85$, soit $85\,\%$.}{on traduit par une phrase}
\end{correction}

\etape{À vous de jouer}
\begin{minipage}[c]{0.6\linewidth}\raggedright
\begin{questions}[label=\alph*)]
  \item Une urne contient $15$ boules indiscernables, numérotées de $1$ à $15$. On en tire
        une au hasard. On note $A$ \og obtenir un multiple de $3$\fg{} et $B$ \og obtenir un nombre
        pair\fg{}. Calculer $P(A)$, $P(B)$ et $P(A\cap B)$.
  \item Le diagramme ci-contre décrit les $120$ visiteurs des portes ouvertes d'un lycée
        ($D$ : démonstration de drones ; $I$ : impression 3D). On choisit un visiteur au
        hasard. Calculer $P(D\cap I)$, $P(D\cup I)$ et $P(\overline{D}\cap\overline{I})$.
\end{questions}
\end{minipage}\hfill
\begin{minipage}[c]{0.36\linewidth}\centering
\begin{tikzpicture}[scale=0.8]
  \venn{D}{I}{$40$}{$30$}{$20$}{$30$}
\end{tikzpicture}
\end{minipage}
\reponse{a) $P(A)=\frac{1}{3}$ ; $P(B)=\frac{7}{15}$ ; $P(A\cap B)=\frac{2}{15}$ ;\;
b)~$0{,}25$ ; $0{,}75$ ; $0{,}25$.}
\end{methodebox}

\afaire{exercices 9 à 12}

\partiecours{Loi de probabilité}

Quand les issues n'ont pas toutes la même chance (dé truqué, relevés sur une longue
période), on donne la probabilité de chacune dans un tableau.

\begin{definitionbox}[Définition 6 --- Loi de probabilité]
Donner la \textbf{loi de probabilité} d'une expérience, c'est donner la probabilité de
chacune de ses issues, le plus souvent dans un tableau.
\end{definitionbox}

\begin{proprietebox}[Propriété 2 --- Somme des probabilités]
\begin{itemize}
  \item La somme des probabilités de toutes les issues est égale à $1$.
  \item La probabilité d'un événement est la somme des probabilités des issues qui le
        réalisent.
\end{itemize}
\end{proprietebox}

\begin{methodebox}[Méthode 4 --- Utiliser une loi de probabilité]
\etapeexemples
Sur une aire de covoiturage, on choisit au hasard une voiture qui part et on compte les
personnes à bord. Le tableau donne la loi de probabilité.
\begin{center}
\renewcommand{\arraystretch}{1.4}
\begin{tabular}{|l|c|c|c|c|}
\hline
Personnes à bord & $2$ & $3$ & $4$ & $5$\\
\hline
Probabilité & $0{,}45$ & $0{,}30$ & $p$ & $0{,}05$\\
\hline
\end{tabular}
\end{center}
\begin{questions}[label=\alph*)]
  \item Calculer $p$.
  \item Calculer la probabilité de l'événement $A$ \og il y a au moins $3$ personnes à bord\fg{}.
  \item Calculer la probabilité de l'événement $B$ \og il y a un nombre pair de personnes à
        bord\fg{}.
\end{questions}
\begin{correction}
\cas{a)}
\begin{calculs}
0{,}45+0{,}30+p+0{,}05 &= 1 &\qquad& \rem{la somme des probabilités vaut $1$}\\
0{,}80+p &= 1 &&\\
p &= 1-0{,}80 &&\\
p &= \res{0{,}20} &&
\end{calculs}
\cas{b)}
\begin{calculs}
P(A) &= 0{,}30+0{,}20+0{,}05 &\qquad& \rem{$A=\ens{3\,;4\,;5}$ : on additionne}\\
P(A) &= \res{0{,}55} &&
\end{calculs}
\cas{c)}
\begin{calculs}
P(B) &= 0{,}45+0{,}20 &\qquad& \rem{$B=\ens{2\,;4}$}\\
P(B) &= \res{0{,}65} &&
\end{calculs}
\end{correction}

\etape{À vous de jouer}
Lors d'un tournoi, un judoka dispute quatre combats. Le tableau donne la loi de
probabilité du nombre de combats qu'il gagne.
\begin{center}
\renewcommand{\arraystretch}{1.4}
\begin{tabular}{|l|c|c|c|c|c|}
\hline
Combats gagnés & $0$ & $1$ & $2$ & $3$ & $4$\\
\hline
Probabilité & $0{,}10$ & $0{,}25$ & $0{,}35$ & $p$ & $0{,}10$\\
\hline
\end{tabular}
\end{center}
\begin{questions}[label=\alph*)]
  \item Calculer $p$.
  \item Calculer la probabilité de l'événement $A$ \og il gagne au moins $2$ combats\fg{}.
  \item Calculer la probabilité de l'événement $B$ \og il gagne un nombre impair de
        combats\fg{}.
\end{questions}
\reponse{a) $p=0{,}20$ ;\; b)~$P(A)=0{,}65$ ;\; c)~$P(B)=0{,}45$.}
\end{methodebox}

\afaire{exercices 13 à 15}

\partiecours{Événement contraire, réunion}

\begin{proprietebox}[Propriété 3 --- Contraire et réunion]
\begin{minipage}[c]{0.62\linewidth}\raggedright
Pour deux événements $A$ et $B$ :
\begin{itemize}
  \item $P(\overline{A})=1-P(A)$ ;
  \item $P(A\cup B)=P(A)+P(B)-P(A\cap B)$ ;
  \item si $A$ et $B$ sont incompatibles, $P(A\cup B)=P(A)+P(B)$.
\end{itemize}
Dans $P(A)+P(B)$, les issues de $A\cap B$ sont comptées deux fois : on retire une fois
$P(A\cap B)$.
\end{minipage}\hfill
\begin{minipage}[c]{0.34\linewidth}\centering
\begin{tikzpicture}[scale=0.7]
  \begin{scope}\clip (0,0) circle (1.1); \fill[coulMeth!30] (1,0) circle (1.1);\end{scope}
  \vennbord
  \node[font=\scriptsize,align=center] at (0.5,-1.8) {$A\cap B$ : compté dans $A$\\et
    dans $B$};
\end{tikzpicture}
\end{minipage}
\end{proprietebox}

\begin{methodebox}[Méthode 5 --- Calculer la probabilité d'un contraire, d'une réunion]
\etapeexemples
Une usine fabrique des panneaux solaires. On choisit un panneau au hasard à la sortie de la
chaîne et on note $A$ \og le panneau a un défaut de cadre\fg{}, $B$ \og le panneau a un défaut de
câblage\fg{}, $C$ \og le panneau a un défaut de verre\fg{}. On sait que $P(A)=0{,}04$ ;
$P(B)=0{,}03$ ; $P(A\cap B)=0{,}01$ ; $P(C)=0{,}02$ ; et aucun panneau n'a à la fois un
défaut de cadre et un défaut de verre.
\begin{questions}[label=\alph*)]
  \item Calculer $P(\overline{A})$ et traduire le résultat par une phrase.
  \item Calculer la probabilité qu'un panneau ait au moins un des défauts $A$ et $B$.
  \item En déduire la probabilité qu'un panneau n'ait ni défaut de cadre ni défaut de
        câblage.
  \item Calculer $P(A\cup C)$.
\end{questions}
\begin{correction}
\cas{a)}
\begin{calculs}
P(\overline{A}) &= 1-P(A) &\qquad& \rem{probabilité du contraire}\\
P(\overline{A}) &= 1-0{,}04 = \res{0{,}96} &&
\end{calculs}
\explique{$96\,\%$ des panneaux n'ont pas de défaut de cadre.}{$\overline{A}$ : \og le
panneau n'a pas de défaut de cadre\fg{}}
\cas{b)}
\begin{calculs}
P(A\cup B) &= P(A)+P(B)-P(A\cap B) &\qquad& \rem{\og au moins un des deux\fg{} : la réunion}\\
P(A\cup B) &= 0{,}04+0{,}03-0{,}01 &&\\
P(A\cup B) &= \res{0{,}06} &&
\end{calculs}
\cas{c)}
\begin{calculs}
P(\overline{A}\cap\overline{B}) &= 1-P(A\cup B) &\qquad& \rem{\og ni l'un ni l'autre\fg{} :
contraire de \og au moins un\fg{}}\\
P(\overline{A}\cap\overline{B}) &= 1-0{,}06 = \res{0{,}94} &&
\end{calculs}
\cas{d)}
\begin{calculs}
P(A\cup C) &= P(A)+P(C) &\qquad& \rem{$A$ et $C$ sont incompatibles}\\
P(A\cup C) &= 0{,}04+0{,}02 = \res{0{,}06} &&
\end{calculs}
\end{correction}

\etape{À vous de jouer}
Une usine fabrique des batteries de trottinette. On choisit une batterie au hasard et on
note $A$ \og la batterie a un défaut d'autonomie\fg{}, $B$ \og la batterie a un défaut de
connecteur\fg{}, $C$ \og la batterie a un défaut de boîtier\fg{}. On sait que $P(A)=0{,}05$ ;
$P(B)=0{,}04$ ; $P(A\cap B)=0{,}02$ ; $P(C)=0{,}03$ ; et aucune batterie n'a à la fois un
défaut d'autonomie et un défaut de boîtier.
\begin{questions}[label=\alph*)]
  \item Calculer $P(\overline{A})$ et traduire le résultat par une phrase.
  \item Calculer la probabilité qu'une batterie ait au moins un des défauts $A$ et $B$.
  \item En déduire la probabilité qu'une batterie n'ait ni défaut d'autonomie ni défaut de
        connecteur.
  \item Calculer $P(A\cup C)$.
\end{questions}
\reponse{a) $P(\overline{A})=0{,}95$ : $95\,\%$ des batteries n'ont pas de défaut
d'autonomie ;\; b)~$0{,}07$ ;\; c)~$0{,}93$ ;\; d)~$0{,}08$.}
\end{methodebox}

\afaire{exercices 16 à 19}

\partiecours{Arbre des possibles}

Une expérience peut se faire en plusieurs étapes : taper les deux chiffres d'un code, tirer
deux jetons l'un après l'autre. Un arbre montre alors toutes les issues.

\begin{definitionbox}[Définition 7 --- Arbre des possibles]
Dans un \textbf{arbre des possibles}, chaque étape de l'expérience forme un niveau de
branches ; chaque chemin, du départ au bout de l'arbre, est une issue.

\smallskip
Quand chaque nœud d'une même étape porte le même nombre de branches, le nombre d'issues est
le \textbf{produit} des nombres de branches de chaque étape.
\end{definitionbox}

\begin{methodebox}[Méthode 6 --- Dénombrer les issues avec un arbre]
\etapeexemples
Un cadenas de vélo s'ouvre avec un code de deux chiffres, chacun égal à $1$, à $2$ ou à $3$.
On essaie un code au hasard.
\begin{questions}[label=\alph*)]
  \item Représenter toutes les issues par un arbre. Combien y en a-t-il ?
  \item Calculer la probabilité de trouver le bon code.
  \item Calculer la probabilité de l'événement $A$ \og les deux chiffres sont différents\fg{}.
\end{questions}
Un sac contient quatre jetons numérotés de $1$ à $4$. On tire un jeton au hasard, on le
garde, puis on en tire un deuxième : le premier donne le chiffre des dizaines, le second
celui des unités.
\begin{questions}[label=\alph*),start=4]
  \item Représenter toutes les issues par un arbre. Combien y en a-t-il ?
  \item Calculer la probabilité de l'événement $B$ \og le nombre formé est supérieur à $40$\fg{}.
  \item Calculer la probabilité de l'événement $C$ \og le nombre formé est un multiple de $3$\fg{}.
\end{questions}
\begin{correction}
\cas{a)}\par\nopagebreak
\begin{minipage}[c]{0.58\linewidth}\raggedright
Premier chiffre : $3$ branches ; au bout de chacune, deuxième chiffre : $3$ branches.
\begin{calculs}
3\times 3 &= \res{9} &\qquad& \rem{$3$ choix, puis $3$ choix}
\end{calculs}
Il y a $9$ codes possibles.
\end{minipage}\hfill
\begin{minipage}[c]{0.38\linewidth}\centering
\begin{tikzpicture}[x=1cm,y=0.4cm,font=\small]
  \node (O) at (0,0) {};
  \foreach \ca/\ya in {1/3,2/0,3/-3}{%
    \node[evt] (N\ca) at (1.2,\ya) {$\ca$};
    \draw (O.center) -- (N\ca);
    \foreach \cb/\db in {1/1,2/0,3/-1}{%
      \node[evt] (N\ca\cb) at (2.4,\ya+\db) {$\cb$};
      \draw (N\ca) -- (N\ca\cb);
      \node[right,font=\footnotesize] at (N\ca\cb.east) {\ca\cb};}}
\end{tikzpicture}
\end{minipage}
\cas{b)}
\explique{Un seul code est le bon, parmi $9$ codes qui ont la même chance : la probabilité
de le trouver est $\res{\dfrac{1}{9}}$.}{code choisi au hasard : équiprobabilité}
\cas{c)}
\begin{calculs}
P(A) &= \dfrac{6}{9} = \res{\dfrac{2}{3}} &\qquad& \rem{$A=\ens{12\,;13\,;21\,;23\,;31\,;32}$ ;
on simplifie par $3$}
\end{calculs}
\cas{d)}\par\nopagebreak
\begin{minipage}[c]{0.58\linewidth}\raggedright
Premier jeton : $4$ branches ; il n'est pas remis dans le sac, donc il reste $3$ jetons
pour le deuxième : $3$ branches au bout de chacune.
\begin{calculs}
4\times 3 &= \res{12} &\qquad& \rem{$4$ choix, puis $3$ choix}
\end{calculs}
Il y a $12$ nombres possibles.
\end{minipage}\hfill
\begin{minipage}[c]{0.38\linewidth}\centering
\begin{tikzpicture}[x=1cm,y=0.4cm,font=\small]
  \node (O) at (0,0) {};
  \foreach \ca/\ya/\reste in {1/4.5/{2,3,4},2/1.5/{1,3,4},3/-1.5/{1,2,4},4/-4.5/{1,2,3}}{%
    \node[evt] (N\ca) at (1.2,\ya) {$\ca$};
    \draw (O.center) -- (N\ca);
    \foreach \cb [count=\k] in \reste {%
      \node[evt] (N\ca\cb) at (2.4,\ya+2-\k) {$\cb$};
      \draw (N\ca) -- (N\ca\cb);
      \node[right,font=\footnotesize] at (N\ca\cb.east) {\ca\cb};}}
\end{tikzpicture}
\end{minipage}
\cas{e)}
\begin{calculs}
P(B) &= \dfrac{3}{12} = \res{\dfrac{1}{4}} &\qquad& \rem{$B=\ens{41\,;42\,;43}$ ; on
simplifie par $3$}
\end{calculs}
\cas{f)}
\begin{calculs}
P(C) &= \dfrac{4}{12} = \res{\dfrac{1}{3}} &\qquad& \rem{$C=\ens{12\,;21\,;24\,;42}$ ; on
simplifie par $4$}
\end{calculs}
\end{correction}

\etape{À vous de jouer}
Un antivol s'ouvre avec un code formé d'une lettre, X ou Y, suivie d'un chiffre, égal à
$1$, à $2$ ou à $3$. On essaie un code au hasard.
\begin{questions}[label=\alph*)]
  \item Représenter toutes les issues par un arbre. Combien y en a-t-il ?
  \item Calculer la probabilité de trouver le bon code.
  \item Calculer la probabilité de l'événement $A$ \og le code se termine par $3$\fg{}.
\end{questions}
Un sac contient trois jetons numérotés de $1$ à $3$. On tire un jeton au hasard, on le
garde, puis on en tire un deuxième : le premier donne le chiffre des dizaines, le second
celui des unités.
\begin{questions}[label=\alph*),start=4]
  \item Représenter toutes les issues par un arbre. Combien y en a-t-il ?
  \item Calculer la probabilité de l'événement $B$ \og le nombre formé est supérieur à $30$\fg{}.
  \item Calculer la probabilité de l'événement $C$ \og le nombre formé est impair\fg{}.
\end{questions}
\reponse{a) $2\times 3=6$ issues ;\; b)~$\frac{1}{6}$ ;\; c)~$\frac{1}{3}$ ;\; d)~$3\times
2=6$ issues ;\; e)~$\frac{1}{3}$ ;\; f)~$\frac{2}{3}$.}
\end{methodebox}

\afaire{exercices 20 à 22}

\partiecours{Probabilité conditionnelle et tableau croisé}

\textbf{But :} on sait déjà qu'un événement $A$ est réalisé (le vélo loué est électrique) et
on cherche la probabilité d'un autre événement $B$ (la location est à la journée). On ne
regarde plus toute la population, mais seulement celle de $A$.

\begin{definitionbox}[Définition 8 --- Probabilité conditionnelle]
Soit $A$ un événement de probabilité non nulle. La \textbf{probabilité de $B$ sachant $A$},
notée $P_A(B)$, est la probabilité que $B$ soit réalisé quand on sait que $A$ l'est.

\smallskip
Dans un tableau d'effectifs, on ne garde que la ligne ou la colonne de $A$ :
\[
  P_A(B)=\dfrac{\text{effectif de $A\cap B$}}{\text{effectif de $A$}}.
\]
\end{definitionbox}

\begin{remarquebox}[À ne pas confondre]
Ces trois probabilités ont le même numérateur, l'effectif de $A\cap B$, mais pas le même
dénominateur :
\begin{itemize}
  \item $P(A\cap B)$ : on divise par l'effectif \textbf{total} (\og $A$ et $B$\fg{}) ;
  \item $P_A(B)$ : on divise par l'effectif de $A$ (\og parmi les $A$, la part des $B$\fg{}) ;
  \item $P_B(A)$ : on divise par l'effectif de $B$ (\og parmi les $B$, la part des $A$\fg{}).
\end{itemize}
Les mots \og sachant que\fg{} et \og parmi\fg{} annoncent une probabilité conditionnelle.
\end{remarquebox}

\begin{methodebox}[Méthode 7 --- Calculer une probabilité conditionnelle avec un tableau croisé]
\etapeexemples
Un loueur de vélos a enregistré $400$ locations cet été. Le tableau, à compléter, les range
selon le vélo et la durée de la location.
\begin{center}
\renewcommand{\arraystretch}{1.4}
\begin{tabular}{|l|c|c|c|}
\hline
 & Électrique & Classique & Total\\
\hline
À la journée & $150$ & $90$ & \\
\hline
À l'heure & $50$ & & \\
\hline
Total & & $200$ & $400$\\
\hline
\end{tabular}
\end{center}
On choisit une location au hasard ; on note $E$ \og le vélo est électrique\fg{} et $J$ \og la
location est à la journée\fg{}.
\begin{questions}[label=\alph*)]
  \item Compléter le tableau.
  \item Calculer $P(E)$ et $P(E\cap J)$.
  \item Calculer $P_E(J)$ et traduire le résultat par une phrase.
  \item Calculer $P_J(E)$ et traduire le résultat par une phrase.
\end{questions}
\begin{correction}
\cas{a)}\par\nopagebreak
\begin{minipage}[c]{0.5\linewidth}\raggedright
\begin{calculs}
150+90 &= \res{240} &\qquad& \rem{total à la journée}\\
150+50 &= \res{200} && \rem{total des électriques}\\
200-90 &= \res{110} && \rem{classiques à l'heure}\\
50+110 &= \res{160} && \rem{total à l'heure}
\end{calculs}
\end{minipage}\hfill
\begin{minipage}[c]{0.46\linewidth}\centering
\renewcommand{\arraystretch}{1.3}
\begin{tabular}{|l|c|c|c|}
\hline
 & Électrique & Classique & Total\\
\hline
À la journée & $150$ & $90$ & $240$\\
\hline
À l'heure & $50$ & $110$ & $160$\\
\hline
Total & $200$ & $200$ & $400$\\
\hline
\end{tabular}
\end{minipage}
\cas{b)}
\begin{calculs}
P(E) &= \dfrac{200}{400} = \res{0{,}5} &\qquad& \rem{$200$ électriques sur $400$ locations}\\
P(E\cap J) &= \dfrac{150}{400} = \res{0{,}375} && \rem{électriques \textbf{et} à la journée,
sur le total}
\end{calculs}
\cas{c)}
\begin{calculs}
P_E(J) &= \dfrac{150}{200} = \res{0{,}75} &\qquad& \rem{on ne garde que la colonne
\og Électrique\fg{}}
\end{calculs}
\explique{$75\,\%$ des locations de vélos électriques sont à la journée.}{\og parmi les vélos
électriques\fg{}}
\cas{d)}
\begin{calculs}
P_J(E) &= \dfrac{150}{240} = \res{0{,}625} &\qquad& \rem{on ne garde que la ligne \og À la
journée\fg{}}
\end{calculs}
\explique{$62{,}5\,\%$ des locations à la journée concernent un vélo électrique.}{$P_E(J)$
et $P_J(E)$ ne sont pas égales}
\end{correction}

\etape{À vous de jouer}
Un club de judo compte $250$ adhérents. Le tableau, à compléter, les range selon l'âge et
la pratique.
\begin{center}
\renewcommand{\arraystretch}{1.4}
\begin{tabular}{|l|c|c|c|}
\hline
 & Enfant & Adulte & Total\\
\hline
Compétition & $60$ & $20$ & \\
\hline
Loisir & $90$ & & \\
\hline
Total & & $100$ & $250$\\
\hline
\end{tabular}
\end{center}
On choisit un adhérent au hasard ; on note $E$ \og l'adhérent est un enfant\fg{} et $C$
\og l'adhérent fait de la compétition\fg{}.
\begin{questions}[label=\alph*)]
  \item Compléter le tableau.
  \item Calculer $P(E)$ et $P(E\cap C)$.
  \item Calculer $P_E(C)$ et traduire le résultat par une phrase.
  \item Calculer $P_C(E)$ et traduire le résultat par une phrase.
\end{questions}
\reponse{a) totaux en compétition $80$, en loisir $170$, des enfants $150$ ; adultes en
loisir $80$ ;\; b)~$P(E)=0{,}6$ et $P(E\cap C)=0{,}24$ ;\;
c)~$P_E(C)=0{,}4$ : $40\,\%$ des enfants font de la compétition ;\; d)~$P_C(E)=0{,}75$ :
$75\,\%$ des compétiteurs sont des enfants.}
\end{methodebox}

\afaire{exercices 23 à 26}

\partiecours{Formule des probabilités conditionnelles}

Souvent, on ne connaît pas les effectifs, seulement des pourcentages. On passe de la
définition 8 à une formule sur les probabilités en divisant le numérateur et le
dénominateur par l'effectif total $N$ :
\begin{calculs}
P_A(B) &= \dfrac{\text{effectif de $A\cap B$}}{\text{effectif de $A$}} &\qquad& \rem{définition 8}\\
P_A(B) &= \dfrac{\dfrac{\text{effectif de $A\cap B$}}{N}}{\dfrac{\text{effectif de $A$}}{N}}
  && \rem{on divise en haut et en bas par $N$}\\
P_A(B) &= \dfrac{P(A\cap B)}{P(A)} && \rem{chaque quotient est une probabilité}
\end{calculs}
Avec les vélos de la partie VII : $\dfrac{P(E\cap J)}{P(E)}=\dfrac{0{,}375}{0{,}5}=0{,}75$,
c'est bien $P_E(J)$.

\begin{proprietebox}[Propriété 4 --- Formule des probabilités conditionnelles]
Si $P(A)\neq 0$ :
\begin{center}
$P_A(B)=\dfrac{P(A\cap B)}{P(A)}$ \qquad et donc \qquad $P(A\cap B)=P(A)\times P_A(B)$.
\end{center}
\end{proprietebox}

\begin{remarquebox}[Traduire un énoncé]
On choisit un salarié d'une entreprise au hasard ; $A$ : \og il a un abonnement de
transport\fg{}, $V$ : \og il a un vélo électrique\fg{}.
\begin{itemize}
  \item \og $30\,\%$ des salariés ont un abonnement\fg{} : $P(A)=0{,}3$ ;
  \item \og $12\,\%$ des salariés ont un abonnement \textbf{et} un vélo électrique\fg{} :
        $P(A\cap V)=0{,}12$ ;
  \item \og \textbf{parmi} les abonnés, $40\,\%$ ont un vélo électrique\fg{}, ou \og $40\,\%$
        \textbf{des abonnés} ont un vélo électrique\fg{} : $P_A(V)=0{,}4$.
\end{itemize}
\end{remarquebox}

\begin{methodebox}[Méthode 8 --- Traduire un énoncé et utiliser la formule]
\etapeexemples
\begin{questions}[label=\alph*)]
  \item Chez un concessionnaire, $60\,\%$ des voitures vendues sont électriques, et
        $45\,\%$ des voitures vendues sont électriques et équipées d'une pompe à chaleur.
        On choisit une vente au hasard ; on note $E$ \og la voiture est électrique\fg{} et $C$
        \og la voiture a une pompe à chaleur\fg{}. Traduire les données par des probabilités,
        puis calculer $P_E(C)$ et traduire le résultat par une phrase.
  \item Dans une salle d'escalade, $30\,\%$ des grimpeurs sont débutants ; parmi les
        débutants, $80\,\%$ louent leurs chaussons. On choisit un grimpeur au hasard ; on
        note $D$ \og le grimpeur est débutant\fg{} et $L$ \og le grimpeur loue ses chaussons\fg{}.
        Traduire les données par des probabilités, puis calculer $P(D\cap L)$ et traduire
        le résultat par une phrase.
\end{questions}
\begin{correction}
\cas{a)}
\explique{Données : $P(E)=\res{0{,}6}$ et $P(E\cap C)=\res{0{,}45}$.}{\og électriques
\textbf{et} équipées\fg{} : une intersection}
\begin{calculs}
P_E(C) &= \dfrac{P(E\cap C)}{P(E)} &\qquad& \rem{formule de la propriété 4}\\
P_E(C) &= \dfrac{0{,}45}{0{,}6} &&\\
P_E(C) &= \res{0{,}75} &&
\end{calculs}
\explique{$75\,\%$ des voitures électriques vendues ont une pompe à chaleur.}{$P_E(C)$ :
parmi les électriques}
\cas{b)}
\explique{Données : $P(D)=\res{0{,}3}$ et $P_D(L)=\res{0{,}8}$.}{\og parmi les débutants\fg{} :
probabilité sachant $D$}
\begin{calculs}
P(D\cap L) &= P(D)\times P_D(L) &\qquad& \rem{formule de la propriété 4}\\
P(D\cap L) &= 0{,}3\times 0{,}8 &&\\
P(D\cap L) &= \res{0{,}24} &&
\end{calculs}
\explique{$24\,\%$ des grimpeurs sont débutants et louent leurs chaussons.}{$D$
\textbf{et} $L$, sur l'ensemble des grimpeurs}
\end{correction}

\etape{À vous de jouer}
\begin{questions}[label=\alph*)]
  \item Dans une compagnie de bus, $70\,\%$ des bus sont électriques, et $56\,\%$ des bus
        sont électriques et climatisés. On choisit un bus au hasard ; on note $E$ \og le bus
        est électrique\fg{} et $C$ \og le bus est climatisé\fg{}. Traduire les données par des
        probabilités, puis calculer $P_E(C)$ et traduire le résultat par une phrase.
  \item Dans un lycée, $40\,\%$ des élèves viennent à vélo ; parmi eux, $25\,\%$ ont un vélo
        à assistance électrique. On choisit un élève au hasard ; on note $V$ \og l'élève vient
        à vélo\fg{} et $A$ \og l'élève a un vélo à assistance électrique\fg{}. Traduire les données
        par des probabilités, puis calculer $P(V\cap A)$ et traduire le résultat par une
        phrase.
\end{questions}
\reponse{a) $P(E)=0{,}7$ et $P(E\cap C)=0{,}56$ ; $P_E(C)=0{,}8$ : $80\,\%$ des bus
électriques sont climatisés ;\; b)~$P(V)=0{,}4$ et $P_V(A)=0{,}25$ ; $P(V\cap A)=0{,}1$ :
$10\,\%$ des élèves viennent avec un vélo à assistance électrique.}
\end{methodebox}

\afaire{exercices 27 à 29}

\partiecours{Arbre pondéré}

Quand une expérience se déroule en deux temps (le temps qu'il fait, puis la livraison), on
la représente par un arbre dont les branches portent des probabilités.

\begin{definitionbox}[Définition 9 --- Arbre pondéré]
\begin{minipage}[c]{0.5\linewidth}\raggedright
Un \textbf{arbre pondéré} porte une probabilité sur chaque branche :
\begin{itemize}
  \item au premier niveau, $P(A)$ et $P(\overline{A})$ ;
  \item au second niveau, des probabilités conditionnelles : après $A$, $P_A(B)$ et
        $P_A(\overline{B})$ ; après $\overline{A}$, $P_{\overline{A}}(B)$ et
        $P_{\overline{A}}(\overline{B})$.
\end{itemize}
Chaque chemin mène à une intersection : $A\cap B$, $A\cap\overline{B}$\dots
\end{minipage}\hfill
\begin{minipage}[c]{0.46\linewidth}\centering
\begin{tikzpicture}
  \arbre{A}{B}{P(A)}{P(\overline{A})}{P_A(B)}{P_A(\overline{B})}%
    {P_{\overline{A}}(B)}{P_{\overline{A}}(\overline{B})}
  \feuilles{A}{B}
\end{tikzpicture}
\end{minipage}
\end{definitionbox}

\begin{proprietebox}[Propriété 5 --- Règles de calcul sur un arbre]
\begin{itemize}
  \item La somme des probabilités des branches issues d'un même nœud est égale à $1$ ;
        par exemple, $P(\overline{A})=1-P(A)$ et $P_A(\overline{B})=1-P_A(B)$.
  \item La probabilité d'un chemin est le \textbf{produit} des probabilités de ses
        branches : $P(A\cap B)=P(A)\times P_A(B)$ (propriété 4).
\end{itemize}
\end{proprietebox}

\begin{methodebox}[Méthode 9 --- Construire un arbre pondéré, calculer la probabilité d'un chemin]
\etapeexemples
Une entreprise livre des colis par drone. $75\,\%$ des livraisons ont lieu par beau temps.
Par beau temps, $96\,\%$ des livraisons arrivent à l'heure ; par mauvais temps, $72\,\%$.
On choisit une livraison au hasard ; on note $B$ \og la livraison a lieu par beau temps\fg{} et
$H$ \og la livraison arrive à l'heure\fg{}.
\begin{questions}[label=\alph*)]
  \item Construire l'arbre pondéré de la situation.
  \item Calculer la probabilité qu'une livraison ait lieu par beau temps et arrive à
        l'heure ; traduire le résultat par une phrase.
  \item Calculer la probabilité qu'une livraison ait lieu par mauvais temps et arrive en
        retard.
\end{questions}
\begin{correction}
\cas{a)}\par\nopagebreak
\begin{minipage}[c]{0.5\linewidth}\raggedright
Données : $P(B)=0{,}75$ ; $P_B(H)=0{,}96$ ; $P_{\overline{B}}(H)=0{,}72$. Les autres
branches s'en déduisent : branches d'un même nœud, somme $1$.
\begin{calculs}
P(\overline{B}) &= 1-0{,}75 = \res{0{,}25} &&\\
P_B(\overline{H}) &= 1-0{,}96 = \res{0{,}04} &&\\
P_{\overline{B}}(\overline{H}) &= 1-0{,}72 = \res{0{,}28} &&
\end{calculs}
\end{minipage}\hfill
\begin{minipage}[c]{0.46\linewidth}\centering
\begin{tikzpicture}
  \arbre{B}{H}{0{,}75}{0{,}25}{0{,}96}{0{,}04}{0{,}72}{0{,}28}
  \feuilles{B}{H}
\end{tikzpicture}
\end{minipage}
\cas{b)}
\begin{calculs}
P(B\cap H) &= P(B)\times P_B(H) &\qquad& \rem{chemin $B$ puis $H$ : on multiplie}\\
P(B\cap H) &= 0{,}75\times 0{,}96 &&\\
P(B\cap H) &= \res{0{,}72} &&
\end{calculs}
\explique{$72\,\%$ des livraisons ont lieu par beau temps et arrivent à l'heure.}{$B$
\textbf{et} $H$, sur l'ensemble des livraisons}
\cas{c)}
\begin{calculs}
P(\overline{B}\cap\overline{H}) &= P(\overline{B})\times P_{\overline{B}}(\overline{H})
  &\qquad& \rem{chemin $\overline{B}$ puis $\overline{H}$}\\
P(\overline{B}\cap\overline{H}) &= 0{,}25\times 0{,}28 &&\\
P(\overline{B}\cap\overline{H}) &= \res{0{,}07} &&
\end{calculs}
\end{correction}

\etape{À vous de jouer}
Une maison est équipée de panneaux solaires. $60\,\%$ des jours de l'année sont
ensoleillés. Un jour ensoleillé, la production couvre toute la consommation de la maison
avec une probabilité de $0{,}9$ ; un jour sans soleil, avec une probabilité de $0{,}45$. On
choisit un jour au hasard ; on note $S$ \og le jour est ensoleillé\fg{} et $C$ \og la production
couvre la consommation\fg{}.
\begin{questions}[label=\alph*)]
  \item Construire l'arbre pondéré de la situation.
  \item Calculer la probabilité qu'un jour soit ensoleillé et que la production couvre la
        consommation ; traduire le résultat par une phrase.
  \item Calculer la probabilité qu'un jour soit sans soleil et que la production ne couvre
        pas la consommation.
\end{questions}
\reponse{a) $P(\overline{S})=0{,}4$ ; $P_S(\overline{C})=0{,}1$ ;
$P_{\overline{S}}(\overline{C})=0{,}55$ ;\; b)~$P(S\cap C)=0{,}54$ ;\;
c)~$P(\overline{S}\cap\overline{C})=0{,}22$.}
\end{methodebox}

\afaire{exercices 30 à 32}

\partiecours{Formule des probabilités totales}

\textbf{But :} un événement peut se réaliser de deux façons ; une livraison arrive à l'heure
par beau temps \emph{ou} par mauvais temps. Pour obtenir sa probabilité, on additionne les
probabilités des chemins qui y mènent.

\begin{proprietebox}[Propriété 6 --- Formule des probabilités totales]
\begin{minipage}[c]{0.5\linewidth}\raggedright
L'événement $B$ se réalise avec $A$ ou avec $\overline{A}$ :
\begin{calculs}
P(B) &= P(A\cap B)+P(\overline{A}\cap B) &&\\
P(B) &= P(A)\times P_A(B)+P(\overline{A})\times P_{\overline{A}}(B) &&
\end{calculs}
Sur l'arbre, $P(B)$ est la somme des probabilités des chemins qui mènent à $B$ (en rouge).
\end{minipage}\hfill
\begin{minipage}[c]{0.46\linewidth}\centering
\begin{tikzpicture}[arbA/.style={chemin},arbAB/.style={chemin},arbnA/.style={chemin},
    arbnAB/.style={chemin}]
  \arbre{A}{B}{P(A)}{P(\overline{A})}{P_A(B)}{P_A(\overline{B})}%
    {P_{\overline{A}}(B)}{P_{\overline{A}}(\overline{B})}
  \feuille{AB}{$A\cap B$}
  \feuille{nAB}{$\overline{A}\cap B$}
\end{tikzpicture}
\end{minipage}
\end{proprietebox}

\begin{methodebox}[Méthode 10 --- Calculer une probabilité avec la formule des probabilités totales]
\etapeexemples
\begin{minipage}[c]{0.5\linewidth}\raggedright
Une entreprise livre des colis par drone ; l'arbre ci-contre décrit la situation ($B$ \og la
livraison a lieu par beau temps\fg{}, $H$ \og la livraison arrive à l'heure\fg{}).
\begin{questions}[label=\alph*)]
  \item Calculer $P(H)$.
  \item Une livraison est arrivée à l'heure. Calculer la probabilité qu'elle ait eu lieu
        par beau temps.
\end{questions}
\end{minipage}\hfill
\begin{minipage}[c]{0.46\linewidth}\centering
\begin{tikzpicture}
  \arbre{B}{H}{0{,}75}{0{,}25}{0{,}96}{0{,}04}{0{,}72}{0{,}28}
\end{tikzpicture}
\end{minipage}

\smallskip
Une usine assemble des trottinettes électriques sur deux lignes : la ligne 1 en produit
$60\,\%$, la ligne 2 le reste. À la sortie, $2\,\%$ des trottinettes de la ligne 1 ont un
réglage à reprendre, contre $5\,\%$ de celles de la ligne 2. On choisit une trottinette au
hasard ; on note $L$ \og elle vient de la ligne 1\fg{} et $R$ \og elle a un réglage à reprendre\fg{}.
\begin{questions}[label=\alph*),start=3]
  \item Construire l'arbre pondéré, puis calculer $P(R)$.
  \item Une trottinette a un réglage à reprendre. Calculer la probabilité qu'elle vienne de
        la ligne 1.
\end{questions}
\begin{correction}
\cas{a)}
\explique{Deux chemins mènent à $H$ : $B$ puis $H$, et $\overline{B}$ puis $H$.}{$H$ se
réalise par beau ou par mauvais temps}
\begin{calculs}
P(H) &= P(B\cap H)+P(\overline{B}\cap H) &\qquad& \rem{formule des probabilités totales}\\
P(H) &= 0{,}75\times 0{,}96+0{,}25\times 0{,}72 && \rem{produit sur chaque chemin}\\
P(H) &= 0{,}72+0{,}18 &&\\
P(H) &= \res{0{,}9} &&
\end{calculs}
\explique{$90\,\%$ des livraisons arrivent à l'heure.}{on traduit par une phrase}
\cas{b)}
\begin{calculs}
P_H(B) &= \dfrac{P(H\cap B)}{P(H)} &\qquad& \rem{on sait que $H$ est réalisé : propriété 4}\\
P_H(B) &= \dfrac{0{,}72}{0{,}9} && \rem{$P(H\cap B)=P(B\cap H)$, calculé en a)}\\
P_H(B) &= \res{0{,}8} &&
\end{calculs}
\explique{$80\,\%$ des livraisons arrivées à l'heure ont eu lieu par beau temps.}{à ne pas
confondre avec $P_B(H)=0{,}96$}
\cas{c)}\par\nopagebreak
\begin{minipage}[c]{0.5\linewidth}\raggedright
Données : $P(L)=0{,}6$ ; $P_L(R)=0{,}02$ ; $P_{\overline{L}}(R)=0{,}05$. Les autres
branches s'en déduisent : $P(\overline{L})=0{,}4$ ; $P_L(\overline{R})=0{,}98$ ;
$P_{\overline{L}}(\overline{R})=0{,}95$.
\end{minipage}\hfill
\begin{minipage}[c]{0.46\linewidth}\centering
\begin{tikzpicture}
  \arbre{L}{R}{0{,}6}{0{,}4}{0{,}02}{0{,}98}{0{,}05}{0{,}95}
\end{tikzpicture}
\end{minipage}
\begin{calculs}
P(R) &= P(L)\times P_L(R)+P(\overline{L})\times P_{\overline{L}}(R) &\qquad& \rem{deux
chemins mènent à $R$}\\
P(R) &= 0{,}6\times 0{,}02+0{,}4\times 0{,}05 &&\\
P(R) &= 0{,}012+0{,}02 &&\\
P(R) &= \res{0{,}032} &&
\end{calculs}
\cas{d)}
\begin{calculs}
P_R(L) &= \dfrac{P(R\cap L)}{P(R)} &\qquad& \rem{on sait que $R$ est réalisé}\\
P_R(L) &= \dfrac{0{,}012}{0{,}032} && \rem{$P(R\cap L)=0{,}6\times 0{,}02$, calculé en c)}\\
P_R(L) &= \res{0{,}375} &&
\end{calculs}
\explique{$37{,}5\,\%$ des trottinettes à régler viennent de la ligne 1.}{la ligne 1
produit plus, mais se trompe moins}
\end{correction}

\etape{À vous de jouer}
\begin{minipage}[c]{0.5\linewidth}\raggedright
Une maison est équipée de panneaux solaires ; l'arbre ci-contre décrit la situation ($S$ \og le
jour est ensoleillé\fg{}, $C$ \og la production couvre la consommation\fg{}).
\begin{questions}[label=\alph*)]
  \item Calculer $P(C)$.
  \item Un jour, la production a couvert la consommation. Calculer la probabilité que ce
        jour ait été ensoleillé.
\end{questions}
\end{minipage}\hfill
\begin{minipage}[c]{0.46\linewidth}\centering
\begin{tikzpicture}
  \arbre{S}{C}{0{,}6}{0{,}4}{0{,}9}{0{,}1}{0{,}45}{0{,}55}
\end{tikzpicture}
\end{minipage}

\smallskip
Lors d'un concours de robotique, $70\,\%$ des robots sont à roues, les autres à chenilles.
$20\,\%$ des robots à roues réussissent l'épreuve tout-terrain, contre $70\,\%$ des robots à
chenilles. On choisit un robot au hasard ; on note $R$ \og le robot est à roues\fg{} et $T$ \og le
robot réussit l'épreuve tout-terrain\fg{}.
\begin{questions}[label=\alph*),start=3]
  \item Construire l'arbre pondéré, puis calculer $P(T)$.
  \item Un robot a réussi l'épreuve tout-terrain. Calculer la probabilité qu'il soit à
        roues.
\end{questions}
\reponse{a) $P(C)=0{,}72$ ;\; b)~$P_C(S)=0{,}75$ ;\; c)~$P(T)=0{,}35$ ;\;
d)~$P_T(R)=0{,}4$.}
\end{methodebox}

\afaire{exercices 33 à 35}

\partiecours{Événements indépendants}

\textbf{But :} savoir si le fait que $A$ soit réalisé change la probabilité de $B$. Par
exemple, la pluie change-t-elle la probabilité qu'un tramway soit à l'heure ?

\begin{definitionbox}[Définition 10 --- Événements indépendants]
Soit $A$ un événement de probabilité non nulle. L'événement $B$ est \textbf{indépendant}
de $A$ si $P_A(B)=P(B)$ : savoir que $A$ est réalisé ne change pas la probabilité de $B$.
\end{definitionbox}

\begin{proprietebox}[Propriété 7 --- Indépendance et produit]
Soit $A$ et $B$ deux événements de probabilités non nulles. Si $B$ est indépendant de $A$,
alors
\begin{center}
$P(A\cap B)=P(A)\times P(B)$ \qquad et \qquad $P_B(A)=P(A)$ :
\end{center}
$A$ est aussi indépendant de $B$. On dit simplement que $A$ et $B$ sont
\textbf{indépendants}.
\end{proprietebox}

\demo{\textbf{But :} on veut démontrer que, si $P_A(B)=P(B)$, alors $P(A\cap B)=P(A)\times
P(B)$ et $P_B(A)=P(A)$.}
\begin{calculs}
P(A\cap B) &= P(A)\times P_A(B) &\qquad& \rem{propriété 4}\\
P(A\cap B) &= P(A)\times P(B) && \rem{car $P_A(B)=P(B)$}
\end{calculs}
\begin{calculs}
P_B(A) &= \dfrac{P(A\cap B)}{P(B)} &\qquad& \rem{propriété 4, en échangeant $A$ et $B$}\\
P_B(A) &= \dfrac{P(A)\times\barre{P(B)}}{\barre{P(B)}} && \rem{ligne précédente ; on
simplifie par $P(B)$}\\
P_B(A) &= P(A) &&
\end{calculs}
\textbf{Conclusion :} $A$ est indépendant de $B$.

\begin{methodebox}[Méthode 11 --- Justifier ou utiliser l'indépendance de deux événements]
\etapeexemples
\begin{questions}[label=\alph*)]
  \item Sur une ligne de tramway, on a relevé $200$ trajets. On choisit un trajet au
        hasard ; on note $R$ \og il pleut pendant le trajet\fg{} et $H$ \og le tramway est à
        l'heure\fg{}. Les événements $R$ et $H$ sont-ils indépendants ?
        \begin{center}
        \renewcommand{\arraystretch}{1.4}
        \begin{tabular}{|l|c|c|c|}
        \hline
         & Pluie & Pas de pluie & Total\\
        \hline
        À l'heure & $45$ & $135$ & $180$\\
        \hline
        En retard & $5$ & $15$ & $20$\\
        \hline
        Total & $50$ & $150$ & $200$\\
        \hline
        \end{tabular}
        \end{center}
  \item Une entreprise livre des colis par drone. Par beau temps, $96\,\%$ des livraisons
        arrivent à l'heure ; au total, $90\,\%$ des livraisons arrivent à l'heure. On note
        $B$ \og la livraison a lieu par beau temps\fg{} et $H$ \og la livraison arrive à
        l'heure\fg{}. Les événements $B$ et $H$ sont-ils indépendants ?
  \item Un circuit comporte deux composants montés en série, qui tombent en panne
        indépendamment l'un de l'autre. Sur un an, le composant 1 fonctionne avec une
        probabilité de $0{,}95$ et le composant 2 avec une probabilité de $0{,}9$. Le circuit
        fonctionne quand les deux composants fonctionnent. Calculer la probabilité que le
        circuit fonctionne pendant un an, puis la probabilité qu'il tombe en panne.
\end{questions}
\begin{correction}
\cas{a)}
\begin{calculs}
P_R(H) &= \dfrac{45}{50} = \res{0{,}9} &\qquad& \rem{colonne \og Pluie\fg{} : $50$ trajets}\\
P(H) &= \dfrac{180}{200} = \res{0{,}9} && \rem{sur l'ensemble des $200$ trajets}
\end{calculs}
\explique{Conclusion : $P_R(H)=P(H)$, donc \res{$R$ et $H$ sont indépendants} : la pluie ne
change pas la probabilité que le tramway soit à l'heure.}{on compare $P_R(H)$ à $P(H)$ ;
le tramway roule sur une voie réservée}
\cas{b)}
\explique{$P_B(H)=0{,}96$ et $P(H)=0{,}9$.}{\og par beau temps\fg{} : sachant $B$ ; \og au
total\fg{} : $P(H)$}
\explique{Conclusion : $P_B(H)\neq P(H)$, donc \res{$B$ et $H$ ne sont pas
indépendants} : le beau temps augmente la probabilité d'arriver à l'heure.}{une différence
suffit}
\cas{c)}
\explique{On note $F_1$ \og le composant 1 fonctionne\fg{} et $F_2$ \og le composant 2
fonctionne\fg{}.}{on nomme les événements}
\begin{calculs}
P(F_1\cap F_2) &= P(F_1)\times P(F_2) &\qquad& \rem{$F_1$ et $F_2$ sont indépendants
(propriété 7)}\\
P(F_1\cap F_2) &= 0{,}95\times 0{,}9 &&\\
P(F_1\cap F_2) &= \res{0{,}855} &&\\
P(\overline{F_1\cap F_2}) &= 1-0{,}855 = \res{0{,}145} && \rem{la panne est le
contraire du fonctionnement}
\end{calculs}
\explique{Le circuit fonctionne pendant un an avec une probabilité de $0{,}855$ ; il tombe
en panne avec une probabilité de $0{,}145$.}{on traduit par une phrase}
\end{correction}

\etape{À vous de jouer}
\begin{questions}[label=\alph*)]
  \item Une station de vélos en libre-service a enregistré $300$ locations. On choisit une
        location au hasard ; on note $W$ \og la location a lieu le week-end\fg{} et $E$ \og le vélo
        est électrique\fg{}. Les événements $W$ et $E$ sont-ils indépendants ?
        \begin{center}
        \renewcommand{\arraystretch}{1.4}
        \begin{tabular}{|l|c|c|c|}
        \hline
         & Week-end & Semaine & Total\\
        \hline
        Électrique & $40$ & $80$ & $120$\\
        \hline
        Classique & $60$ & $120$ & $180$\\
        \hline
        Total & $100$ & $200$ & $300$\\
        \hline
        \end{tabular}
        \end{center}
  \item Une maison est équipée de panneaux solaires. Un jour ensoleillé, la production
        couvre la consommation avec une probabilité de $0{,}9$ ; sur l'ensemble des jours de
        l'année, avec une probabilité de $0{,}72$. On note $S$ \og le jour est ensoleillé\fg{} et
        $C$ \og la production couvre la consommation\fg{}. Les événements $S$ et $C$ sont-ils
        indépendants ?
  \item Un vélo électrique roule en mode assisté quand sa batterie et son moteur
        fonctionnent ; ils tombent en panne indépendamment l'un de l'autre. Sur un an, la
        batterie fonctionne avec une probabilité de $0{,}9$ et le moteur avec une probabilité
        de $0{,}8$. Calculer la probabilité que l'assistance fonctionne pendant un an, puis la
        probabilité qu'elle tombe en panne.
\end{questions}
\reponse{a) $P_W(E)=0{,}4=P(E)$ : indépendants ;\; b)~$P_S(C)=0{,}9\neq 0{,}72$ : non
indépendants ;\; c)~$0{,}72$ puis $0{,}28$.}
\end{methodebox}

\afaire{exercices 36 à 38}

\end{document}
